\subsection{Sandboxing y Seguridad}

La Sección 2.2 comparó las plataformas disponibles para ejecutar código
generado por LLMs. Esta sección complementa esa discusión enfocándose en el
\textbf{modelo de amenazas} específico de la ejecución de código de IA y las
estrategias de mitigación implementadas en esta investigación.

\subsubsection{Modelo de amenazas}

El código generado por un LLM es inherentemente no confiable. A diferencia del
código escrito por desarrolladores humanos ---que pasa por procesos de revisión,
testing y control de versiones---, el código generado por un modelo puede
contener errores sutiles producto de alucinaciones, efectos no intencionados de
prompt injection, o simplemente lógica incorrecta que el modelo no tiene
capacidad de verificar antes de la ejecución. Esto establece un modelo de
amenazas con cuatro categorías principales:

\begin{enumerate}
    \item \textbf{Escape del sandbox}: código que explota vulnerabilidades del
    entorno de ejecución para obtener acceso al sistema host. En el caso de
    Docker, esto incluye exploits del kernel compartido, escape mediante
    \texttt{/proc} o \texttt{/sys}, y abuso de capabilities del contenedor.

    \item \textbf{Exfiltración de datos}: código que utiliza acceso a red para
    enviar datos sensibles a servidores externos. Un LLM que genera código con
    llamadas HTTP ocultas podría exfiltrar datos del volumen montado o
    información del entorno de ejecución.

    \item \textbf{Denegación de servicio}: código que consume recursos del host
    de manera desproporcionada ---bucles infinitos, asignaciones masivas de
    memoria, fork bombs--- degradando o interrumpiendo el servicio para otros
    procesos en la máquina.

    \item \textbf{Acceso al sistema host}: código que intenta acceder al
    sistema de archivos del host fuera de los volúmenes explícitamente montados,
    escalar privilegios, o manipular procesos del host a través de interfaces
    compartidas.
\end{enumerate}

Es importante notar que estas amenazas no requieren intención maliciosa del
usuario. Un LLM puede generar código peligroso como efecto secundario de
instrucciones benignas ---por ejemplo, un modelo que intenta ``optimizar''
una consulta podría generar código que itera indefinidamente, o un modelo que
intenta ``explorar'' el sistema podría generar código que accede a rutas
sensibles.

\subsubsection{Estrategias de mitigación en Docker}

La implementación de esta tesis utiliza Docker como entorno de ejecución con
una configuración de seguridad que mitiga cada categoría de amenaza
identificada. Los contenedores se ejecutan con los siguientes flags de
seguridad:

\begin{lstlisting}[language=bash,caption={Flags de seguridad Docker utilizados en esta investigación}]
docker run \
  --network none \
  --memory 512m \
  --cpus 1 \
  --read-only \
  --cap-drop ALL \
  --security-opt no-new-privileges \
  -v ./data:/data:rw
\end{lstlisting}

Cada flag aborda amenazas específicas:

\begin{itemize}
    \item \textbf{\texttt{--network none}}: elimina completamente el acceso a
    red del contenedor, mitigando la exfiltración de datos. El código
    ejecutado no puede establecer conexiones salientes ni recibir conexiones
    entrantes.

    \item \textbf{\texttt{--memory 512m}}: establece un límite máximo de
    memoria de 512 megabytes. Si el proceso excede este límite, el kernel
    del host termina el contenedor mediante OOM killer, mitigando la
    denegación de servicio por consumo de memoria.

    \item \textbf{\texttt{--cpus 1}}: limita el contenedor a un núcleo de CPU,
    previniendo que un proceso intensivo monopolice los recursos de cómputo
    del host.

    \item \textbf{\texttt{--read-only}}: monta el sistema de archivos raíz del
    contenedor como solo lectura. El código ejecutado no puede modificar
    binarios del sistema, instalar software adicional ni alterar la
    configuración del contenedor. Únicamente el directorio \texttt{/data},
    montado explícitamente como lectura-escritura, es modificable.

    \item \textbf{\texttt{--cap-drop ALL}}: elimina todas las Linux
    capabilities del proceso del contenedor. Las capabilities son permisos
    granulares que el kernel otorga a procesos ---como \texttt{NET\_RAW}
    (sockets raw), \texttt{SYS\_ADMIN} (operaciones de administración) o
    \texttt{MKNOD} (creación de dispositivos). Al eliminarlas todas, el
    proceso ejecuta con el mínimo privilegio posible.

    \item \textbf{\texttt{--security-opt no-new-privileges}}: previene que el
    proceso o sus hijos adquieran privilegios adicionales mediante mecanismos
    como \texttt{setuid}, \texttt{setgid} o ejecución de binarios con
    capabilities de archivo. Esto cierra vectores de escalamiento de
    privilegios dentro del contenedor.
\end{itemize}

Adicionalmente, los contenedores son \textbf{efímeros}: se crean para cada
ejecución de código y se destruyen inmediatamente después. Esto garantiza que
no existe persistencia de estado entre ejecuciones, eliminando la posibilidad
de que código malicioso deje artefactos que afecten ejecuciones posteriores.

\subsubsection{Alternativas de aislamiento}

Aunque Docker proporciona un nivel de aislamiento adecuado para esta
investigación, existen alternativas con garantías de seguridad más fuertes:

\paragraph{gVisor.} Desarrollado por Google \citep{gvisor2024}, gVisor
implementa un kernel de aplicación en espacio de usuario que intercepta todas
las llamadas al sistema (syscalls) del contenedor. A diferencia de Docker
estándar, donde los contenedores comparten el kernel del host y las syscalls se
ejecutan directamente, gVisor actúa como un proxy que reimplementa un
subconjunto del kernel Linux en Go. Esto reduce drásticamente la superficie de
ataque: un exploit de kernel en un contenedor con gVisor afecta únicamente al
kernel de aplicación de gVisor, no al kernel del host. El trade-off es un
overhead de rendimiento en operaciones intensivas de syscalls.

\paragraph{Firecracker.} Como se describió en la Sección 2.2, Firecracker
\citep{e2b2024sandbox} proporciona aislamiento a nivel de microVM, ejecutando
un kernel independiente para cada sandbox. Esto elimina por completo el vector
de escape de kernel compartido, ya que cada ejecución opera sobre su propio
kernel. El costo es mayor consumo de recursos por instancia comparado con
contenedores.

\paragraph{WebAssembly.} El sandboxing a nivel de instrucción de WebAssembly,
descrito en la Sección 2.2, ofrece aislamiento por construcción: el código
ejecuta dentro de un módulo binario que no tiene acceso directo al sistema
operativo. Sin embargo, el soporte limitado de bibliotecas Python (a través de
Pyodide) lo hace menos adecuado para escenarios que requieren dependencias
científicas o de análisis de datos.

El trade-off general entre estas alternativas es claro: \textbf{mayor
aislamiento implica mayor overhead}, lo que potencialmente impacta la latencia
medida en los experimentos. La Tabla~\ref{tab:sandbox-comparison} en la Sección
2.2 cuantifica estas diferencias.

\subsubsection{Implicaciones para la investigación}

La configuración de seguridad tiene implicaciones directas para el diseño
experimental de esta tesis:

\begin{itemize}
    \item \textbf{Variable controlada}: las cuatro configuraciones
    experimentales (A1--A4) utilizan exactamente los mismos flags de seguridad
    de Docker. Esto asegura que cualquier diferencia observada en latencia o
    rendimiento se debe al modo de ejecución (tool calling vs.\ código) y no a
    diferencias en la configuración del sandbox.

    \item \textbf{Componente de la latencia medida}: el overhead de crear un
    contenedor Docker, ejecutar el código y destruir el contenedor es parte de
    la latencia total medida en las configuraciones A2 y A4 (ejecución de
    código). Este overhead es un costo inherente del paradigma de ejecución de
    código que los resultados deben reflejar de manera transparente.

    \item \textbf{Limitación del timeout}: cada ejecución tiene un timeout
    configurable (30 segundos por defecto) que actúa como una capa adicional
    de protección contra denegación de servicio y como una variable controlada
    del experimento.
\end{itemize}
